\chapter{Research Methodology}
\label{ch:methodology}

\section{Scope and Approach}

This report is a synthesis document, not a primary research study. It draws on three categories of source material to construct an integrated argument about AI-driven organisational transformation.

\subsection{Peer-Reviewed and Pre-Print Research}

The empirical foundation rests primarily on the work of Dell'Acqua, Mollick, and their collaborators at Harvard Business School. The jagged frontier study \parencite{DellAcqua2026} is a preregistered field experiment with 758 participants, published in \emph{Organization Science} --- a top-tier management journal. The vigilance study \parencite{DellAcqua2023} is a working paper that has been widely cited but has not yet completed peer review. The collaboration modes study \parencite{Randazzo2024} is a working paper with 244 participants.

The task-based automation framework of \textcite{Acemoglu2019} provides the economic theory connecting AI capability to labour market effects. The psychological safety research of \textcite{EdmondsonLei2025} is grounded in the established literature on team effectiveness.

\textbf{Limitation:} The jagged frontier study used GPT-4, now two model generations behind frontier models. The specific task boundaries have likely shifted, though the principle of an uneven frontier remains sound.

\subsection{Consultancy and Industry Reports}

The industry landscape analysis draws on reports from McKinsey \parencite{McKinseyAgenticOrg2026, McKinseySixShifts2026}, BCG \parencite{BCGAgenticShift2025, BCGAgenticEnterprise2025, BCGAIKPIs2025}, Deloitte \parencite{Deloitte2026}, Gartner \parencite{Gartner2025}, the World Economic Forum \parencite{WEF2025}, Capgemini \parencite{Capgemini2025}, and IBM \parencite{IBM2025}.

\textbf{Limitation:} Consultancy reports are produced by organisations with commercial interests in the phenomena they describe. Methodology sections are often absent or sparse. Sample composition and response biases are rarely disclosed. Statistics from these sources are used directionally --- to establish trends and orders of magnitude --- rather than as precise measurements.

\subsection{Practitioner Accounts}

The analysis of Block's intelligence layer model draws on the published paper by \textcite{DorseyBotha2026}, contemporaneous reporting (New York Times, Fortune, The Guardian, VentureBeat), and commentary from \textcite{Zamost2026} and \textcite{Spicer2026}.

The analysis of Every's agent colony model draws on \textcite{Shipper2026} and \textcite{Klaassen2025}, both published on every.to. These are first-person practitioner accounts rather than controlled research.

The analysis of the Dynamic Agentic Mesh draws on operational experience with the VisionClaw system and a codebase audit conducted in April 2026. Claims about the system's capabilities have been reconciled against the code, with corrections noted in Chapter~\ref{ch:technical-substrate}.

\textbf{Limitation:} Practitioner accounts are inherently subjective. The VisionClaw analysis is conducted by the system's architect, introducing an obvious conflict of interest. This conflict is acknowledged in Chapters~\ref{ch:open-questions} and \ref{ch:conclusion}.

\section{Analytical Framework}

The report's analytical framework proceeds in four stages:

\begin{enumerate}
  \item \textbf{Historical analysis}: Establishing hierarchy as an information routing protocol by tracing the structural logic from the Roman army through modern organisational design.
  \item \textbf{Empirical synthesis}: Integrating findings from the peer-reviewed literature, consultancy reports, and practitioner accounts into a coherent picture of the current transformation.
  \item \textbf{Model comparison}: Evaluating three competing organisational models against the empirical evidence, using the jagged frontier and the vigilance problem as analytical lenses.
  \item \textbf{Architectural proposal}: Developing the Dynamic Agentic Mesh as a specific response to the evidence, with explicit acknowledgement of what the evidence supports and what it does not.
\end{enumerate}

\section{Citation Standards}

Where a claim is supported by peer-reviewed evidence, the citation is provided inline. Where a claim is drawn from consultancy reports, the report is cited with a note about its methodological limitations. Where a claim is drawn from practitioner accounts, the source is identified and the anecdotal nature of the evidence is noted.

Where this report makes claims that are not supported by external evidence --- particularly the proposed KPIs, the compounding flywheel model, and the judgment broker role --- this is stated explicitly. Chapter~\ref{ch:open-questions} systematically identifies every known gap in the evidence base.

\section{Disclosure}

The author of this report, Dr.\ John O'Hare, is the founder and principal of DreamLab AI Consulting Ltd and the architect of the VisionClaw system described in Chapter~\ref{ch:technical-substrate}. The report simultaneously analyses the problem space, proposes a solution framework, and describes a commercial offering from the author's firm. This positioning conflict is acknowledged in Chapter~\ref{ch:open-questions} and is offered as context for the reader's own evaluation of the argument.
